% Vietnamese translation for OI Public Library
% Source-PDF: IOI中国国家候选队论文/2002/戴德承.pdf
\documentclass[11pt]{article}
\usepackage{oipl-vn}

\newcommand{\Oh}{\mathcal{O}}
\newcommand{\Area}{\operatorname{Area}}
\newcommand{\sgnarea}{\operatorname{S}}

\title{Lùi một bước, trời cao biển rộng}
\author{Dai Decheng\\Trường Trung học số 1 Sán Đầu}
\date{}

\begin{document}
\maketitle

\origtitle{Lùi một bước, trời cao biển rộng: một vài ứng dụng của tư tưởng cắt bù và chuyển hóa mục tiêu}
\sourcepath{IOI China candidate team papers / 2002 / Dai Decheng.pdf}

\section*{Từ khóa}

Giả thiết nhượng bộ; phương pháp cắt bù; ứng dụng.

\section*{Tóm tắt}

Bài viết này chủ yếu bàn về cách trong thiết kế thuật toán ta chuyển một bài
toán khó giải thành một bài toán có phạm vi lớn hơn nhưng dễ giải hơn, sau đó
cắt bỏ phần đã mở rộng để thu được đáp án cần tìm. Phần thảo luận tập trung
phân tích một số bài toán thường gặp trong hình học tính toán và toán tổ hợp,
đồng thời cũng chạm tới một số nguyên lý cơ bản như nguyên lý bao hàm--loại
trừ. Trong một số bài toán, mục tiêu không phải lúc nào cũng là một thuật toán
hoàn toàn lý tưởng; thay vào đó, ta kết hợp ưu điểm của các phương pháp, đồng
thời bàn cả lợi hại của thuật toán và tác dụng của nó trong quá trình giải
bài.

\section{Khái quát}

Người xưa thường nói: ``lùi một bước, trời cao biển rộng''. Câu ấy dạy ta khi
đối nhân xử thế phải biết nhường nhịn, và khi theo đuổi mục tiêu phải biết
chọn bỏ. Điều này đúng, nhưng ngoài hành vi ứng xử, ta còn có thể rút ra điều
gì từ đó hay không? Chẳng hạn, khi nghiên cứu bị bế tắc, liệu có thể mở rộng
bài toán chưa giải được rồi giải một bài toán rộng hơn không? Câu trả lời là
có. Việc chứng minh giả thuyết Goldbach trong toán học là một ví dụ: ``$1+1$''
chưa làm được thì tạm lùi một bước, thậm chí lùi tới ``$9+9$'', rồi từng bước
tiến gần mục tiêu, cho đến kết quả ``$1+2$'' của Trần Cảnh Nhuận. Dù cuối cùng
giả thuyết ấy có được chứng minh hay không, nó ít nhất cho ta một gợi ý: trong
tin học cũng có thể có sự ``nhượng bộ''; lùi một bước vẫn có thể mở ra đường
rộng.

Khi thiết kế thuật toán, ta thường gặp những bài toán khó giải trực tiếp. Một
cách xử lý là trước hết mở rộng phạm vi cần giải, chẳng hạn nới lỏng điều
kiện, đưa vào giả thiết hoặc phỏng đoán thích hợp, hay phóng đại mục tiêu.
Sau đó ta giải bài toán mới, cuối cùng ``cắt'' phần đã mở rộng đi để giải được
bài toán ban đầu.

Ví dụ, xét bài ``phân tích số nguyên'' ở vòng phân khu NOI 2002, bảng nâng
cao. Bài toán hỏi: phân một số tự nhiên $n$ thành $k$ phần
\[
  A_1,A_2,\ldots,A_k,\qquad A_1\le A_2\le\cdots\le A_k,
\]
có bao nhiêu cách phân khác nhau? Bài toán không thuận tiện để giải trực tiếp,
nên trước hết ta chuyển hóa nó. Không khó thấy bài toán tương đương với việc
đếm số cách phân $n$ thành một số phần tùy ý, trong đó phần lớn nhất bằng
$k$. Gọi $G(n,k)$ là số cách như vậy.

Tuy nhiên, $G(n,k)$ vẫn chưa thuận tiện. Ta đưa vào hàm rộng hơn $F(n,k)$:
số cách phân số tự nhiên $n$ thành một số phần tùy ý, trong đó phần lớn nhất
không vượt quá $k$. Khi đó
\[
\begin{aligned}
  F(n,k)
    &= \sum_{i=1}^{k} G(n,i)
     = \sum_{i=1}^{k} F(n-i,i) \\
    &= F(n-k,k)+F(n,k-1).
\end{aligned}
\]

Như vậy có thể dùng quy hoạch động rất thông thường để giải. Mỗi trạng thái
chỉ cần một phép cộng, độ phức tạp là $\Oh(NK)$, và kết quả bài toán ban đầu
là
\[
  F(n,k)-F(n,k-1).
\]
Ở đây, ban đầu ta chuyển một bài khó giải thành một bài có thể giải nhưng vẫn
cần cộng dồn khá nhiều. Sau đó ta dùng chính ý tưởng ``nhượng bộ'' nói trên:
mở rộng mục tiêu thành ``phần lớn nhất được phép không vượt quá $k$'', cuối
cùng thu được thuật toán bậc hai. Dù vậy, đây chỉ là một bước lùi nhỏ, chủ yếu
để đơn giản hóa thuật toán; nó chưa làm thay đổi căn bản mô hình giải bài.

\section{Ứng dụng trong hình học tính toán}

Hình học tính toán là một lĩnh vực thường xuất hiện trong các kỳ thi tin học
những năm gần đây. Mức độ không quá sâu, nhưng một số vấn đề thường gặp vẫn
làm nhiều thí sinh lúng túng. Có lẽ nguyên nhân là thế giới thực và thế giới
số thuần túy khác nhau khá xa. Ví dụ bài toán che khuất: xác định một vật có
bị một hay vài vật khác che hay không. Trong đời sống hằng ngày, ta chỉ cần
nhìn là biết, nhưng trong hình học tính toán thì không đơn giản như vậy.

Chính vì thế, ``giả thiết nhượng bộ'' có nhiều đất dụng võ trong hình học tính
toán. Nhiều bài toán liên quan đến diện tích đa giác đều dùng tư tưởng này:
tính diện tích đa giác, tính trọng tâm đa giác, xác định quan hệ giữa điểm và
đa giác, v.v. Ngay trong tính toán thường ngày ta cũng đã dùng cách nghĩ ấy.
Chẳng hạn khi tính diện tích phần tô màu bằng diện tích hình chữ nhật trừ diện
tích hình tròn, ta có thể hiểu quá trình như sau: trước hết giả sử hình tròn
cũng là một phần của vùng tô màu, nên toàn bộ là một hình chữ nhật; sau đó lấy
phần ``giả sử'' ấy ra, ta được diện tích cần tìm. Từ tư tưởng này sinh ra một
phương pháp tính diện tích khoa học hơn.

\subsection{Tính diện tích đa giác bất kỳ}

Học hình học tính toán thì không thể tránh bài toán tính diện tích hình phẳng.
Nếu phân theo hình dạng, có thể chia thành hình biên gấp khúc và hình biên
cong; ở đây ta chủ yếu bàn về diện tích hình biên gấp khúc. Hình biên gấp khúc
lại chia thành đa giác lồi và đa giác lõm, giả sử hai cạnh bất kỳ của đa giác
không cắt nhau ngoài các đầu mút.

Diện tích đa giác lồi khá dễ tính: phân nó thành một số tam giác, rồi dùng công
thức Heron hoặc công thức tích có hướng. Với đa giác lõm thì phải suy nghĩ
nhiều hơn. Từ ví dụ cắt bù ở trên, người mới học thường nghĩ tới cách: bù đa
giác lõm thành một đa giác lồi, tính diện tích đa giác lồi rồi trừ phần thừa.
Nhưng cách này có khuyết điểm rõ ràng. Có hình phần thừa chỉ là một tam giác,
có hình phần thừa là một đa giác lồi, và cũng có hình phần thừa lại là một đa
giác lõm. Điều đó trái với nguyên tắc quan trọng nhất của cắt bù: phải thuận
tiện. Cốt lõi của cắt bù là biến phức tạp thành đơn giản; nếu không giảm được
độ phức tạp tính toán, phương pháp ấy mất ý nghĩa.

Thật ra, tính diện tích đa giác bất kỳ đã có một thuật toán kinh điển: công
thức diện tích bằng tích có hướng. Giả sử trong hệ tọa độ Descartes có đa giác
bất kỳ
\[
  A_1A_2\cdots A_k,
\]
với
\[
  A_i=(X_i,Y_i),\qquad 1\le i\le k.
\]
Lấy một điểm tùy ý $O=(X_0,Y_0)$, và hiểu chỉ số $i+1$ theo modulo $k$. Diện
tích có hướng của đa giác là
\[
  \sgnarea =
  \sum_{i=1}^{k}
  \frac12
  \begin{vmatrix}
    X_i-X_0 & X_{i+1}-X_0\\
    Y_i-Y_0 & Y_{i+1}-Y_0
  \end{vmatrix}.
\]
Diện tích hình học là $|\sgnarea|$ nếu các đỉnh được liệt kê theo một chiều
nhất quán.

Trong phương pháp này, diện tích các tam giác được cộng với dấu dương hoặc âm.
Khi $O$ nằm trong đa giác, mọi tam giác đều có cùng dấu và cách tính giống
việc chia đa giác thành các tam giác. Khi $O$ nằm ngoài, những phần dư tự động
bị trừ đi bằng dấu âm, đúng với tinh thần cắt bù.

Tính đúng đắn khá trực quan. Với một đa giác sáu đỉnh như hình trong bản gốc,
nếu ký hiệu $S_{i,j,k,\ldots}$ là diện tích dương của đa giác tạo bởi các điểm
$i,j,k,\ldots$, ta có dạng biến đổi điển hình
\[
\begin{aligned}
  \sgnarea
    &= S_{1,2,0}+S_{2,3,0}+S_{3,4,0}+S_{4,5,0}
       -S_{5,6,0}-S_{6,1,0}\\
    &= S_{1,2,0}+S_{2,3,0}+S_{3,4,0}+S_{4,5,0}
       -S_{5,6,0}-S_{6,1,0}\\
    &\quad +(S_{0,1,3}-S_{0,1,3})
       +(S_{0,1,4}-S_{0,1,4})
       +(S_{0,1,5}-S_{0,1,5})\\
    &= (S_{1,2,0}+S_{2,3,0}-S_{3,1,0})
       +(S_{3,1,0}+S_{3,4,0}-S_{4,1,0})\\
    &\quad +(S_{0,1,4}+S_{0,4,5}-S_{0,1,5})
       +(S_{0,1,5}-S_{0,6,5}-S_{0,1,6})\\
    &= S_{1,2,3}+S_{1,3,4}+S_{1,4,5}+S_{1,5,6}\\
    &= S_{1,2,3,4,5,6}.
\end{aligned}
\]
Đây chỉ chứng minh cho trường hợp minh họa trong hình, nhưng từ đó có thể dễ
dàng suy rộng để chứng minh công thức tổng quát.

Trong ứng dụng thực tế, để đơn giản hóa công thức, ta thường chọn $O$ là gốc
tọa độ. Khi đó
\[
  \sgnarea =
  \sum_{i=1}^{k}
  \frac12
  \begin{vmatrix}
    X_i & X_{i+1}\\
    Y_i & Y_{i+1}
  \end{vmatrix}.
\]
Nếu vẫn chưa hài lòng, ta cũng có thể thử thay điểm tham chiếu bằng một đường
thẳng: chuyển diện tích do điểm và đoạn thẳng quét ra thành diện tích do đoạn
thẳng và đường thẳng tham chiếu quét ra. Bản gốc minh họa cách này bằng hình
vẽ. Dùng lập luận tương tự cũng chứng minh được nó đúng, nhưng so với công
thức dùng điểm tham chiếu thì kém thuận tiện hơn: công thức phức tạp hơn, và
đường thẳng tham chiếu phải chọn sao cho không cắt hình, tốt nhất lại là đường
nằm ngang. Điều này làm giảm đáng kể phạm vi áp dụng. Xét về bản chất, hai
cách cùng một gốc, và đều có độ phức tạp $\Oh(N)$ với $N$ là số đỉnh của đa
giác.

Đây là một ứng dụng cơ bản và thường gặp nhất của phương pháp cắt bù. Nhưng
sức mạnh của nó không chỉ dừng ở đây.

\subsection{Tính trọng tâm đa giác bất kỳ}

``Trọng tâm'' là một khái niệm xuất hiện rất nhiều trong vật lý, và cũng có
không ít người làm tin học thích khai thác nó. Vậy làm thế nào để tính trọng
tâm của một vật?

Một cách từng được giới thiệu là lấy ``chia tam giác'' làm tư tưởng chính:
trước hết phân đa giác thành một số tam giác, lần lượt tính trọng tâm và khối
lượng, tức diện tích, của từng tam giác, rồi hợp nhất các trọng tâm ấy. Bước
chia tam giác là phần phức tạp nhất của thuật toán; độ phức tạp của riêng bước
này khoảng $\Oh(N^2)$, khiến toàn bộ thuật toán vừa phức tạp vừa khó cài đặt.

Ta biết rằng nếu hai chất điểm có khối lượng lần lượt là $M_1,M_2$ và hoành độ
là $X_1,X_2$, thì hoành độ trọng tâm hợp nhất là
\[
  X=\frac{X_1+\lambda X_2}{1+\lambda},
  \qquad
  \lambda=\frac{M_2}{M_1},
  \qquad
  M=M_1+M_2.
\]
Tung độ cũng được tính tương tự. Có thể gọi đây là ``quy tắc cộng trọng tâm'':
$(M_1,X_1)$ và $(M_2,X_2)$ là hai số hạng, còn $(M,X)$ là tổng. Theo định nghĩa
của phép trừ, nếu biết tổng và một số hạng thì ta có thể tìm số hạng còn lại.
Vì vậy, trọng tâm đã có quy tắc cộng thì cũng có quy tắc trừ; dựa vào quan hệ
giữa phép trừ và số âm, ta còn có thể đưa vào khái niệm ``trọng tâm âm'' để
hỗ trợ giải bài.

Chẳng hạn, biết trọng tâm ở $(1,0)$, khối lượng $10\,\mathrm{kg}$, và một phần
có trọng tâm $(0,0)$, khối lượng $5\,\mathrm{kg}$, thì phần còn lại có trọng
tâm $(2,0)$ và khối lượng $5\,\mathrm{kg}$. Như vậy phép toán trọng tâm có
cộng, có trừ, có dương, có âm; do đó có thể dùng cắt bù.

Chuyển công thức diện tích ở phần trước sang bài toán trọng tâm, ta có công
thức sau. Với đa giác có các đỉnh $(X_i,Y_i)$ và khối lượng phân bố đều, chọn
điểm tham chiếu $O=(X_0,Y_0)$. Với mỗi cạnh $A_iA_{i+1}$, đặt
\[
\begin{cases}
  M_i =
  \dfrac12
  \begin{vmatrix}
    X_i-X_0 & X_{i+1}-X_0\\
    Y_i-Y_0 & Y_{i+1}-Y_0
  \end{vmatrix},\\[1.1em]
  P_i=\dfrac{X_0+X_i+X_{i+1}}{3},\\[0.8em]
  Q_i=\dfrac{Y_0+Y_i+Y_{i+1}}{3}.
\end{cases}
\]
Ở đây $M_i$ là diện tích có hướng của tam giác $OA_iA_{i+1}$, còn
$(P_i,Q_i)$ là trọng tâm của tam giác ấy. Nếu chọn $O$ là gốc tọa độ, công
thức $P_i,Q_i$ đúng với dạng khi lấy gốc tọa độ trong bản gốc:
$P_i=(X_i+X_{i+1})/3$ và $Q_i=(Y_i+Y_{i+1})/3$.

Có thể cộng dồn bằng quy tắc trọng tâm:
\[
\begin{cases}
  M\leftarrow M_i,\\
  P\leftarrow P_i,\\
  Q\leftarrow Q_i,
\end{cases}
\qquad\text{khi }M=0,
\]
hoặc nếu $M\ne0$ thì
\[
\begin{cases}
  \lambda=\dfrac{M_i}{M},\\[0.8em]
  P\leftarrow\dfrac{P+\lambda P_i}{1+\lambda},\\[1em]
  Q\leftarrow\dfrac{Q+\lambda Q_i}{1+\lambda},\\[1em]
  M\leftarrow M+M_i.
\end{cases}
\]
Tương đương, ta có thể viết gọn hơn:
\[
  M=\sum_i M_i,\qquad
  P=\frac{\sum_i M_iP_i}{M},\qquad
  Q=\frac{\sum_i M_iQ_i}{M}.
\]
Cuối cùng, $(P,Q)$ là tọa độ trọng tâm cần tìm, còn $M$ là diện tích có hướng
của đa giác. Cách làm này khéo léo dùng khối lượng âm để tận dụng sự gọn gàng
của cắt bù. Mỗi đỉnh chỉ cần xử lý hằng số lần, mỗi lần chỉ vài phép nhân
chia đơn giản, nên độ phức tạp chỉ là $\Oh(N)$; đây là một phương pháp khá lý
tưởng.

\subsection{Xác định quan hệ giữa điểm và đa giác bất kỳ}

Xác định vị trí tương đối giữa điểm và đa giác là một bài toán hình học tính
toán kinh điển, và đã có nhiều lời giải tốt. Trong bài toán này, cắt bù không
có ưu thế rõ rệt; nó chỉ được nêu ra như một cách tham khảo.

Ví dụ, để xác định điểm $O$ có nằm trong hình $K$ hay không, nếu bản gốc vẽ
$K$ như một tứ giác $ABCD$ bị khoét một tam giác $AED$, ta có thể làm như
sau: trước hết kiểm tra $O$ có nằm trong tứ giác $ABCD$ không. Nếu không, hiển
nhiên $O$ không nằm trong $K$. Nếu có, tiếp tục kiểm tra $O$ có nằm trong tam
giác $AED$ không. Nếu có thì $O$ nằm ngoài $K$; nếu không thì $O$ nằm trong
$K$.

Dựa trên tư tưởng ấy, với một số hình phức tạp có thể thiết kế một thuật toán
tổng quát hơn:

\paragraph{\texttt{Inside\_Check}$(P_1,P_2,\ldots,P_N)$.}
\begin{enumerate}
  \item Tìm một đỉnh $P_k$ lồi ra ngoài của hình $K$.
  \item Đặt $h\leftarrow k$.
  \item Trong khi $P_{h+1}$ vẫn lồi ra ngoài và $h+1\ne k$, đặt
  $h\leftarrow h+1$.
  \item Đặt $t\leftarrow k$.
  \item Trong khi $P_{t-1}$ vẫn lồi ra ngoài và $t\ne h$, đặt
  $t\leftarrow t-1$.
  \item Đặt
  \[
    \mathit{flag}\leftarrow
    \texttt{Inside\_Check}(P_t,P_{t+1},\ldots,P_k,P_{k+1},\ldots,P_{h-1},P_h).
  \]
  \item Đặt
  \[
    \mathit{map}\leftarrow
    \texttt{Inside\_Check}(P_h,P_{h+1},\ldots,P_{t-1},P_t).
  \]
  \item Nếu $\mathit{flag}$ đúng và $\mathit{map}$ sai thì trả về
  \textit{true}; ngược lại trả về \textit{false}.
\end{enumerate}

Quá trình trên giải đệ quy. Mỗi lần nó tách hình ban đầu thành một đa giác lồi
chuẩn và một hình không đều, rồi lần lượt kiểm tra điểm có nằm trong hai phần
ấy hay không, cuối cùng suy ra điểm có nằm trong đa giác $K$ hay không.

Tuy nhiên, thuật toán này có khuyết điểm. Bản gốc đưa ra một hình phản ví dụ
trong đó cách tách theo dãy đỉnh lồi không xử lý tốt. Vì vậy có một phiên bản
cải tiến, dùng bao lồi để tránh trường hợp ``vượt khung'':

\paragraph{\texttt{Inside\_Check\_PRO}$(P_1,P_2,\ldots,P_N)$.}
\begin{enumerate}
  \item Tìm bao lồi $P_{a_1},P_{a_2},\ldots,P_{a_k}$ của hình
  $P_1,P_2,\ldots,P_N$.
  \item Đặt
  \[
    \mathit{flag}\leftarrow
    \texttt{Inside}(P_{a_1},P_{a_2},\ldots,P_{a_k}).
  \]
  \item Với mỗi dãy điểm không thuộc bao lồi, cùng với hai đầu mút thuộc bao
  lồi tạo thành một hình $(A_1,A_2,\ldots,A_p)$, cập nhật
  \[
    \mathit{map}\leftarrow
    \mathit{map}\land
    \texttt{Inside\_Check\_PRO}(A_1,A_2,\ldots,A_p).
  \]
  \item Nếu $\mathit{flag}$ đúng và $\mathit{map}$ sai thì trả về
  \textit{true}; ngược lại trả về \textit{false}.
\end{enumerate}

Trong đó \texttt{Inside(P1,P2,\ldots,PN)} cho biết $O$ có nằm trong đa giác
lồi $P_1P_2\cdots P_N$ hay không. Sau khi đưa bao lồi vào, vấn đề của thuật
toán ban đầu được khắc phục. Xét toàn bộ quá trình, mỗi vòng cần tìm một bao
lồi, độ phức tạp $\Oh(N)$ với giả thiết thứ tự biên đã có; cần kiểm tra quan
hệ điểm--đa giác lồi, cũng $\Oh(N)$; và mỗi lần đệ quy ít nhất xóa được một
cạnh khỏi hình ban đầu. Vì số lần đệ quy không vượt quá $N$, độ phức tạp vào
khoảng $\Oh(N^2)$, vẫn có thể chấp nhận.

Dù vậy, ở đây phương pháp cắt bù không thật sự phát huy đặc điểm biến phức
tạp thành đơn giản. Nó phải dùng cắt bù lặp lại nhiều tầng, làm tăng độ phức
tạp thuật toán và độ khó lập trình. So với nó, vài thuật toán $\Oh(N)$ lại
tỏ ra ưu việt hơn. Một cách ngắn gọn là: nếu điểm $O$ nằm trong hình $K$, thì
một tia bất kỳ phát ra từ $O$ sẽ cắt biên của $K$ một số lẻ lần; ngược lại,
nếu $O$ nằm ngoài thì số giao điểm là chẵn. Phương pháp này có độ phức tạp
$\Oh(N)$ và rất dễ cài đặt, nên là một lời giải lý tưởng cho bài toán.

Tóm lại, phương pháp cắt bù sinh ra trong hình học tính toán và cũng được dùng
rất rộng ở đây, từ chứng minh mệnh đề đến các bài tính diện tích, trọng tâm.
Hễ bài toán liên quan đến diện tích hình học, ta thường có thể tìm thấy bóng
dáng của nó. Nhưng qua ví dụ cuối cùng, có thể rút ra một tiền đề khi dùng cắt
bù: phần được bù thêm phải hoàn chỉnh và dễ giải, phần bị cắt đi cũng phải đều
đặn hoặc dễ giải. Nếu không, phương pháp sẽ mất ý nghĩa căn bản. Ứng dụng của
cắt bù trong hình học có thể xem là minh chứng rõ nhất cho ``giả thiết nhượng
bộ'' trong tin học, nhưng giả thiết và cắt bù không chỉ tồn tại trong hình học;
chúng cũng xuất hiện trong nhiều bài toán khác.

\section{Ứng dụng trong toán tổ hợp}

Trong toán tổ hợp, phương pháp cắt bù của hình học tính toán thể hiện thành
``cắt bù'' trong đếm. Khi cần đếm số phần tử trong một phạm vi thỏa điều kiện
nào đó, ta có thể mở rộng phạm vi rồi đếm, chẳng hạn nguyên lý trừ. Hoặc trong
khi thống kê, ta đếm từng phần rồi loại bỏ phần thừa, chẳng hạn nguyên lý bao
hàm--loại trừ. Tóm lại, lùi một bước, trời cao biển rộng: nới điều kiện trước
rồi giải bài thường thuận tiện hơn nhiều.

\subsection{Nguyên lý trừ}

Đây chỉ là một bài toán toán học đơn giản, có thể xem như phần mở rộng của
nguyên lý cộng và nguyên lý nhân. Giả sử từ địa điểm $A$ đến các địa điểm
$B,C,D$ mỗi nơi có $5$ con đường, nhưng đến $E$ chỉ có $3$ con đường. Từ mỗi
nơi $B,C,D,E$ đến $F$ đều có $3$ con đường. Ta có thể tạm giả sử $A$ đến $E$
cũng có $5$ con đường. Khi đó tổng số đường từ $A$ đến $F$ là
\[
  5\times 4\times 3=60.
\]
Trong số đó có
\[
  2\times 3=6
\]
đường do ta ``bịa'' ra, nên số đường thực tế từ $A$ đến $F$ là
\[
  60-6=54.
\]

\subsection{Hoán vị có vùng cấm}

Trước hết ta giới thiệu khái niệm đa thức bàn cờ. Gọi $C$ là một bàn cờ, và
$R_k(C)$ là số cách đặt $k$ quân cờ giống nhau lên $C$. Quy tắc đặt là: sau
khi đặt một quân vào một ô của $C$, hàng và cột chứa ô ấy không được đặt quân
nào khác. Với mọi bàn cờ $C$, quy ước $R_0(C)=1$.

Từ định nghĩa có thể thấy ngay vài ví dụ nhỏ: một ô đơn có $R_1=1$; hai ô độc
lập theo hàng hoặc cột có $R_1=2$; nếu hai ô cùng hàng hoặc cùng cột thì
$R_2=0$; còn hai ô không chung hàng, không chung cột thì $R_2=1$. Các hình vẽ
nhỏ trong bản gốc chỉ minh họa các trường hợp này.

Số cách đặt cờ $R_k(C)$ có các tính chất sau.
\begin{enumerate}
  \item Với mọi bàn cờ $C$ và số nguyên dương $k$, nếu $k$ lớn hơn số ô của
  $C$ thì $R_k(C)=0$.
  \item $R_1(C)$ bằng số ô của $C$.
  \item Nếu bàn cờ $C_1$ qua phép quay hoặc lật trở thành $C_2$, thì
  $R_k(C_1)=R_k(C_2)$.
  \item Lấy một ô chỉ định. Gọi $C_i$ là bàn cờ còn lại sau khi xóa hàng và
  cột chứa ô ấy, còn $C_\ell$ là bàn cờ còn lại sau khi chỉ xóa ô ấy. Khi đó
  \[
    R_k(C)=R_{k-1}(C_i)+R_k(C_\ell),\qquad k\ge1.
  \]
  \item Nếu bàn cờ $C$ gồm hai bàn cờ con $C_1,C_2$ độc lập với nhau, thì
  \[
    R_k(C)=\sum_{i=0}^{k} R_i(C_1)R_{k-i}(C_2).
  \]
\end{enumerate}

\paragraph{Định nghĩa 1.}
Với bàn cờ $C$, đặt
\[
  R(C)=\sum_{k=0}^{\infty} R_k(C)x^k.
\]
Đa thức này được gọi là \term{đa thức bàn cờ}.

Vì $R_k(C)=0$ khi $k$ lớn hơn số ô của bàn cờ, $R(C)$ thực ra chỉ có hữu hạn
hạng. Chẳng hạn, bàn cờ gồm hai ô không chung hàng, không chung cột có đa thức
\[
  R(C)=R_0(C)+R_1(C)x+R_2(C)x^2=1+2x+x^2.
\]

Từ các tính chất của $R_k(C)$ có thể suy ra các tính chất tương ứng của đa
thức bàn cờ.
\begin{enumerate}
  \item Với một ô chỉ định và hai bàn cờ $C_i,C_\ell$ định nghĩa như trên,
  \[
    R(C)=xR(C_i)+R(C_\ell).
  \]
  \item Nếu $C$ được tách thành hai phần độc lập $C_1,C_2$, thì
  \[
    R(C)=R(C_1)R(C_2).
  \]
\end{enumerate}

Dùng hai tính chất này có thể tính đa thức bàn cờ. Ví dụ, với bàn cờ ba ô
hình chữ L, tức một hình vuông $2\times2$ bị bỏ một góc,
\[
\begin{aligned}
  R(C) &= x(1+x)+(1+2x)\\
       &= 1+3x+x^2.
\end{aligned}
\]

Bây giờ ta dùng đa thức bàn cờ để giải bài toán hoán vị có vùng cấm. Một hoán
vị của $X=\{1,2,3,\ldots,n\}$ tương ứng đúng với một cách đặt $n$ quân cờ lên
bàn cờ $n\times n$ sao cho không có hai quân cùng hàng hoặc cùng cột. Ta dùng
$n$ hàng để biểu diễn các phần tử của $X$, các cột để biểu diễn vị trí trong
hoán vị. Ví dụ, cách đặt quân trong hình của bản gốc tương ứng với hoán vị
$2143$. Nếu trong hoán vị yêu cầu phần tử $i$ không được đặt ở vị trí $j$, thì
ô hàng $i$, cột $j$ của bàn cờ không được đặt quân. Tập tất cả các ô cấm ấy
được gọi là \term{vùng cấm}.

\paragraph{Định lý 1.}
Cho bàn cờ $n\times n$ có vùng cấm đã cho, tương ứng với các vị trí không được
phép xuất hiện của các phần tử trong hoán vị của $\{1,2,\ldots,n\}$. Khi đó số
hoán vị thỏa điều kiện vùng cấm là
\[
  n!-r_1(n-1)!+r_2(n-2)!-\cdots+(-1)^nr_n,
\]
trong đó $r_i$ là số cách đặt $i$ quân cờ vào vùng cấm.

\paragraph{Chứng minh.}
Trước hết bỏ qua hạn chế vùng cấm. Số cách đặt $n$ quân không phân biệt lên
bàn cờ $n\times n$, không cùng hàng hoặc cột, là $n!$. Nếu đánh số các quân
là $1,2,\ldots,n$ và xem các cách gán nhãn khác nhau trên cùng tập ô là khác
nhau, số cách đặt quân có nhãn là $n!\cdot n!$. Gọi tập các cách đặt có nhãn
ấy là $S$.

Với $j=1,2,\ldots,n$, gọi $P_j$ là tính chất ``quân cờ số $j$ rơi vào vùng
cấm'', và gọi $A_j$ là tập con của $S$ gồm các cách đặt có tính chất $P_j$.
Ta cần đếm các cách đặt không có quân nào rơi vào vùng cấm.

Số cách để quân số $1$ rơi vào vùng cấm là $r_1$. Sau khi nó đã rơi vào một ô
cấm, các quân $2,3,\ldots,n$ có thể tùy ý đặt trên bàn cờ
$(n-1)\times(n-1)$. Theo nguyên lý nhân,
\[
  |A_1| = r_1 (n-1)!(n-1)!.
\]
Tương tự với mọi $i$, nên
\[
  \sum_{i=1}^{n}|A_i| = r_1(n-1)!\,n!.
\]

Với hai quân số $1$ và $2$ cùng rơi vào vùng cấm, có $2r_2$ cách chọn và gán
hai quân vào hai ô cấm không xung đột. Sau đó các quân còn lại đặt tùy ý trên
bàn cờ $(n-2)\times(n-2)$, do đó với mỗi cặp $i<j$,
\[
  |A_i\cap A_j|=2r_2(n-2)!(n-2)!.
\]
Lấy tổng trên mọi cặp,
\[
\begin{aligned}
  \sum_{1\le i<j\le n}|A_i\cap A_j|
    &= 2\binom{n}{2}r_2(n-2)!(n-2)!\\
    &= r_2(n-2)!\,n!.
\end{aligned}
\]

Tương tự, với giao của $k$ tập $A_i$, tổng trên mọi bộ $k$ chỉ số bằng
\[
  r_k(n-k)!\,n!.
\]
Theo nguyên lý bao hàm--loại trừ, số cách đặt $n$ quân có nhãn mà không quân
nào rơi vào vùng cấm là
\[
  n!\cdot n!-r_1(n-1)!\,n!+r_2(n-2)!\,n!-\cdots+(-1)^nr_n n!.
\]
Mỗi hoán vị không nhãn tương ứng với đúng $n!$ cách gán nhãn cho quân, nên
chia biểu thức trên cho $n!$ ta được công thức của định lý.

Cần nói thêm rằng định lý này thích hợp cho bài toán đặt quân trên bàn cờ
$n\times n$ với vùng cấm tương đối nhỏ. Nếu là bàn cờ $m\times n$ hoặc vùng
cấm rất lớn, thường phải trực tiếp dùng đa thức $R(C)$ để giải.

\paragraph{Ví dụ 1.}
Dùng bốn màu đỏ, lam, lục, vàng để tô bốn thiết bị $A,B,C,D$. Mỗi thiết bị chỉ
dùng một màu và hai thiết bị bất kỳ không được cùng màu. Nếu $B$ không được
dùng màu lam và đỏ, $C$ không được dùng màu lam và lục, $D$ không được dùng
màu lục và vàng, hỏi có bao nhiêu cách tô?

Bài toán này chính là bài đặt cờ có vùng cấm. Đa thức bàn cờ của vùng cấm là
\[
  R(C)=1+6x+10x^2+4x^3.
\]
Suy ra $r_1=6$, $r_2=10$, $r_3=4$. Theo định lý, số phương án là
\[
  N=4!-6\cdot3!+10\cdot2!-4\cdot1!
   =24-36+20-4=4.
\]

\paragraph{Ví dụ 2.}
Bài toán hoán vị sai vị trí cũng có thể xem là bài hoán vị có vùng cấm, trong
đó vùng cấm nằm trên đường chéo chính. Dùng định lý để tính $D_n$.

Vùng cấm gồm $n$ ô độc lập trên đường chéo chính, nên đa thức bàn cờ là
\[
\begin{aligned}
  R(C)
    &= \underbrace{(1+x)(1+x)\cdots(1+x)}_{n\text{ lần}}\\
    &= (1+x)^n\\
    &= 1+\binom{n}{1}x+\binom{n}{2}x^2+\cdots+\binom{n}{n}x^n.
\end{aligned}
\]
Do đó $r_k=\binom{n}{k}$. Thay vào định lý,
\[
\begin{aligned}
  D_n
    &= n!-n(n-1)!+\binom{n}{2}(n-2)!-\cdots
       +(-1)^n\binom{n}{n}0!\\
    &= n!\left(1-\frac1{1!}+\frac1{2!}-\cdots
       +(-1)^n\frac1{n!}\right).
\end{aligned}
\]

\section{Tổng kết}

Qua các ví dụ trên có thể thấy, trong nhiều bài toán khó giải trực tiếp,
``giả thiết nhượng bộ'' và ``phương pháp cắt bù'' có sức mạnh đáng kể. Theo
một nghĩa nào đó, chúng đều phóng đại mục tiêu cần tìm để làm bài toán đơn
giản hơn. Tư tưởng chuyển hóa mô hình này nên trở thành một lối suy nghĩ khá
phổ biến.

Ngoài việc phóng đại mục tiêu, ta còn có thể ``đảo trắng'' mục tiêu, tức là
đi tìm phần bù của tập mục tiêu. Cách nghĩ này cũng đem lại hiệu quả cao trong
nhiều bài toán, với các ứng dụng điển hình như tìm tập độc lập lớn nhất trong
lý thuyết đồ thị, hoặc một số bài toán đếm trong toán tổ hợp.

Bên cạnh đó, bài viết còn đề cập đến nhiều vấn đề chuyển hóa mô hình khác; ở
mỗi ví dụ, chúng đều phát huy tác dụng không thể thay thế. Chuyển hóa mô hình
là một điểm nóng trong các kỳ thi tin học, đồng thời thật sự là một mảng khó
kiểm tra nền tảng toán học. Hai tư tưởng chuyển hóa tiêu biểu được nói tới ở
đây có lẽ chỉ là một phần rất nhỏ. Hy vọng bài viết có thể gợi mở thêm cho
người đọc.

\end{document}
